% Preamble
\documentclass[12pt]{report}
% Packages
\usepackage{amsmath}
\usepackage[margin=20mm]{geometry}
\usepackage{xcolor}
\usepackage[colorlinks=true]{hyperref}
\usepackage{listings}
\usepackage{iftex}
\usepackage{datetime}
\usepackage{longtable}

\setcounter{tocdepth}{1}

\usepackage{mdframed}
\usepackage{xcolor}

\newcommand{\texenginefull}{%
    \ifPDFTeX
    pdf\TeX~\the\pdftexversion.\pdftexrevision
    \fi
    \ifXeTeX
    \XeTeX~\the\XeTeXversion
    \fi
    \ifLuaTeX
    \LuaTeX~\the\luatexversion
    \fi
}

\newmdenv[
    linecolor=red,
    innertopmargin=1em,
    innerbottommargin=1em,
    topline=false,
    bottomline=false,
    rightline=false,
    linewidth=3pt
]{infobox}

\newcommand{\step}[1]{\textbf{Step #1:}\quad}
% Metadata
\title{Administrator Guide\\Version 0.1}
\author{Leo LO}
\date{5 Feb 2026}
% Document
\begin{document}
    \maketitle
    \chapter*{Version History}
    \begin{tabular}{|l|l|l|p{10cm}|}
        \hline
        Version & Date & Author & Notes \\\hline
        0.1 & Feb 2026 & Leo Lo & Initial Release \\\hline
    \end{tabular}
    \tableofcontents

    \pagebreak
    \hspace{0pt}
    \vfill
    The key words ``MUST'', ``MUST NOT'', ``REQUIRED'', ``SHALL'', ``SHALL NOT'', ``SHOULD'', ``SHOULD NOT'',
    ``RECOMMENDED'', ``NOT RECOMMENDED'', ``MAY'', and ``OPTIONAL'' in this document are to be interpreted as described
    in \href{https://datatracker.ietf.org/doc/html/bcp14}{BCP 14}\footnote{https://datatracker.ietf.org/doc/html/bcp14}
    [\href{https://datatracker.ietf.org/doc/html/rfc2119}{RFC2119}]\footnote{https://datatracker.ietf.org/doc/html/rfc2119}
    [\href{https://datatracker.ietf.org/doc/html/rfc8174}{RFC8174}]\footnote{https://datatracker.ietf.org/doc/html/rfc8174}
    when, and only when, they appear in all capitals, as shown here.
    \vfill
    \hspace{0pt}
    Build by \texenginefull\ on \today\ at \currenttime
    \vfill
    \hspace{0pt}
    \pagebreak
    \chapter{Common Administrator Task}\label{ch:common-tasks}
    \section{All password lost}\label{sec:all-password-lost}
    If all the administrator's password are lost, and need to be reset, this cannot be done in the system itself, but
    via updating the password directly.

    \begin{infobox}
        This steps can also update regular user's password.
        But you SHOULD follow the usual step to reset the password as this method will require the password to be
        shared via plaintext.
    \end{infobox}

    In order to reset the password in this manner, you MUST have access to the following things ---
    \begin{itemize}
        \item Access to API call \texttt{/api/util/password}.
        This is allowed by default by the system but may not be allowed based on the web server settings.
        \item Access to the database, and able execute ``UPDATE'' to the database table ``user''.
    \end{itemize}

    \begin{infobox}
        It is RECOMMENDED to use a randomly generated password
    \end{infobox}

    \step{1} Make an API call to \texttt{/api/util/password}, this API call will hash the password for you.
    An example will be \texttt{/api/util/password?password=YOUR\_PASSWORD\_HERE}.
    
    \step{2} Note the password hash returned from the API call.
    It is located at a JSON property ``password''.

    \step{3} Execute and commit the following SQL on the database.
    Replace the value as instructed
    \begin{lstlisting}[language=SQL,
           basicstyle=\ttfamily,
           numbers=left,
           numberstyle=\tiny
    ]
UPDATE `user` SET
    password = '<password hash>',
    last_password_date = NOW(),
    force_password_change = 1
WHERE username = '<username>';
    \end{lstlisting}

    \begin{infobox}
        You MAY skip the `force\_password\_change` changes if you do not wish you force a password changes afterwards.
        This is only recommended if the account is being used by another automated system for security reason.

        Please note that the password may be stored as plaintext in the server access log, depending on the settings.
    \end{infobox}

    \chapter{Configurations}\label{ch:configurations}
    \section{Configuration Infomations}\label{sec:configuration-infomations}
    \begin{longtable}{|p{.40\textwidth}|p{.15\textwidth}|p{.45\textwidth}|}\hline
    \multicolumn{3}{|c|}{Authentication}\\\hline
    Configuration Name & Type & Description \\\hline\endhead
    allow\_register & boolean & True if users are allowed to register themselves. \\\hline
    auth.force\_pwd\_chg\_token\_age & Integer & The number of seconds which a password reset token is valid for
    when it is generated for a forced password change.
    This number MUST be greater than 0, and SHOULD be long enough for a user to set their password.
    It is RECOMMENDED to set the value no less than 300 seconds, but no more than 3600 seconds for security reason.
    \\\hline
    auth.permission\_cache & Integer & The number of seconds which the user's role will be cached.
    This cache can be manually pruned by administrator if needed.\\\hline
    auth.token\_life & Integer & The number which a authentication token is valid  \\\hline
    authfail.factor\_base & Double & The base amount of the delay when there are authentication failure.
    This is the value $\mbox{Delay}_{base}$ in section~\ref{sec:response-delay}\\\hline
    authfail.factor\_exp & Double & The exponent amount of the delay when there are authentication failure.
    This is the value $\mbox{exponent}$ in section~\ref{sec:response-delay}.
    This value SHOULD be greater than 1. \\\hline
    authfail.min\_delay & Integer & The minimum amount of the delay when there are authentication failure.
    This is the value $\mbox{Delay}_{min}$ in section~\ref{sec:response-delay}. \\\hline
    authfail.max\_delay& Integer &  The maximum amount of the delay when there are authentication failure.
    This is the value $\mbox{Delay}_{min}$ in section~\ref{sec:response-delay}.
    Thus value MUST be greater than the configuration for authfail.min\_delay. \\\hline
    \end{longtable}
    \begin{longtable}{|p{.40\textwidth}|p{.15\textwidth}|p{.45\textwidth}|}\hline
    \multicolumn{3}{|c|}{IRC}\\\hline
    Configuration Name & Type & Description \\\hline\endhead
        irc.enabled & Boolean & Enable the IRC bot function.\\\hline
        irc.networks & List & A comma separated list of network keys.
        This will be part of the name of network specific configuration in this section and will be denoted as
        ``network'' in this document\\\hline
    \end{longtable}
    \begin{longtable}{|p{.40\textwidth}|p{.15\textwidth}|p{.45\textwidth}|}\hline
    \multicolumn{3}{|c|}{Job}\\\hline
    Configuration Name & Type & Description \\\hline\endhead
        job.log\_data & boolean & Log data received in the job log if the job support this feature.

        \textbf{WARNING:} Although some jobs support this, they may not work if the data is too long.
        The database is NOT designed to log a large amount of data.\\\hline
    job.threadpool.size & Integer & Number of thread to process the queueing job. \\\hline
    system\_job\_runner & String & The username of the user to own the job started by a timer \\\hline
    \multicolumn{3}{|l|}{Please refer to job detail page for job specific configurations}\\\hline
    \end{longtable}
    \begin{longtable}{|p{.40\textwidth}|p{.15\textwidth}|p{.45\textwidth}|}\hline
    \multicolumn{3}{|c|}{Network Rail Data Stream}\\\hline
    Configuration Name & Type & Description \\\hline\endhead
    network.stomp.enabled & Boolean & Is the data stream from Network Rail enabled \\\hline
    networkrail.batch\_size & Integer & The batch size when importing data \\\hline
    networkrail.username & String & The username to login to Network Rail data \\\hline
    networkrail.password & Password & The password to login to Network Rail data \\\hline
    networkrail.schedule\_cache & Integer & Number of days which the UID - schedule mapping will be created \\\hline
    schedule.retain\_period & Integer & Number of days which the past data will be retained.
    Unless the schedule had been stared by an user.\\\hline
    \multicolumn{3}{|l|}{Please refer to section~\ref{subsec:network-rail-data-stream} on how to register}\\\hline
    \end{longtable}
    \begin{longtable}{|p{.40\textwidth}|p{.15\textwidth}|p{.45\textwidth}|}\hline
    \multicolumn{3}{|c|}{Schedule}\\\hline
    Configuration Name & Type & Description \\\hline\endhead
    schedule.default\_size & Integer & The default size of the result when searching for schedules \\\hline
    schedule.max\_size & Integer & The maximum size of the result when searching for schedules \\\hline
    \end{longtable}


    \section{Registration of external feed}\label{sec:registration-of-external-feed}
    \subsection{Network Rail}\label{subsec:network-rail-data-stream}
    You may register to Network Rail data stream at \href{https://publicdatafeeds.networkrail.co.uk}{https://publicdatafeeds.networkrail.co.uk}.
    Then you can enter the username and password in the configuration.

    \begin{infobox}
        The password will be stored in plaintext in the system, and may be transmitted to remote system in plaintext.
        You SHOULD use a unique password for this system.
    \end{infobox}

    \chapter{Implementation details}\label{ch:implementation-details}
    This section will provide some implementation detail which may affect how configuration are set.
    \section{Response Delay}\label{sec:response-delay}
    When authentication failed, the delay of n-th failure will be calculated as follows:

    $$\mbox{Delay}_{base} = \mbox{Base Time}\times\mbox{exponent}^{n}$$
    $$\mbox{Delay}_{final} = clamp(\mbox{Delay}_{base},\mbox{Delay}_{min},\mbox{Delay}_{max})\footnote{%
        \texttt{clamp(value, lower, upper)} ensures \texttt{value} stays within $[\text{lower}, \text{upper}]$. Equivalently, $\min(\text{upper}, \max(\text{value}, \text{lower}))$.%
    }$$

    \begin{infobox}
        Although it is permitted to have exponent be less than 1, it is NOT RECOMMENDED as this will reduce the amount
        of delay as the number of failure increases.
        This is contra to usual security protocol as the delay should be increases when there are more failures.
    \end{infobox}

\end{document}